% Breadth values are generated from
% benchmark_results/stage_c_reproducibility_manifest.json; factorial values
% come from the matched C1xC2/C1xC2xC4 reports and the post-crash all-15 C4
% decision artifact.
% Regenerate the manifest with:
% serverTest/.venv/bin/python serverTest/build_stage_c_reproducibility_manifest.py

\section{Long-Horizon Component Ablations}
\label{sec:stage-c-long-horizon}

We evaluate the maintained Stage-C components with a matched long-horizon
protocol: 24 landmark clusters, 30 outer iterations, one nominal local
nonlinear step, shared-only full block-metric camera consensus, direct
left-$\mathrm{SE}(3)$ tangent normal equations, points-$p_{95}$ scene
normalization, initial Jacobi camera scaling, persistent DABA trust-region
state, and Schur--PCG with relative tolerance $10^{-2}$.  This retained
plain/C1/C5 matrix is distinct from the post-crash C4 carry-over gate below,
which uses curvature recovery and finite Nesterov as its named reference.  All methods minimize
the standard Snavely pixel sum-squared error (SSE).  Ratios below one favor the
candidate.  Worker CPU is the sum of user and system time reported by
\texttt{/usr/bin/time}; optimization time is coordinator-reported elapsed time.

\subsection{Breadth of safeguarded acceleration and adaptive local work}

Table~\ref{tab:stage-c-breadth} reports the two components with complete
long-horizon breadth.  C1 is safeguarded Themelis-style outer acceleration with
a nominal DRS fallback and restart.  C5 adapts each cluster between one and two
local nonlinear steps using the measured proximal defect.  C5 does not add
proximal oracle calls; it changes the amount of work within a nominal oracle.

\begin{table}[t]
  \centering
  \caption{Matched K24/I30 Stage-C breadth.  SSE, optimization, oracle, and CPU
  columns are geometric-mean ratios to the plain reference.  W/T/L counts
  per-scene endpoint SSE.}
  \label{tab:stage-c-breadth}
  \small
  \begin{tabular}{lrrrrrr}
    \hline
    Method & SSE & W/T/L & Worst & Optimization & Prox. oracles & Worker CPU \\
    \hline
    \multicolumn{7}{l}{\emph{SfM\_Init-derived 1DSfM, 15 scenes}} \\
    C1 & 0.850595 & 13/0/2 & 1.045075 & 1.541091 & 1.673801 & 1.726258 \\
    C5 & 0.952889 & 13/0/2 & 1.016898 & 1.176137 & 1.000000 & 1.468372 \\
    C1+C5 & \textbf{0.786739} & \textbf{15/0/0} & \textbf{0.947483}
      & 1.832873 & 1.667622 & 2.544050 \\
    \hline
    \multicolumn{7}{l}{\emph{Large BAL, 3 scenes}} \\
    C1 & 0.991375 & 2/0/1 & 1.002193 & 1.724159 & 1.700226 & 1.753382 \\
    C5 & 0.999815 & 1/2/0 & 1.000000 & 1.114341 & 1.000000 & 1.198071 \\
    C1+C5 & 0.996026 & 2/0/1 & 1.016363 & 1.913371 & 1.633218 & 2.043677 \\
    \hline
  \end{tabular}
\end{table}

C1 retains its short-horizon quality gains on 1DSfM: it improves 13 of 15
scenes and reduces geometric-mean SSE by $14.9\%$.  Its large-BAL endpoint gain
is only $0.86\%$, while worker CPU increases by $75.3\%$; C1 is therefore a
quality mechanism rather than a universal speedup.  Isolated C5 improves 13 of
15 1DSfM scenes without adding proximal calls.  C5 activates depth two in
$10.9\%$ of nominal cluster-iterations.  It is inactive on BAL1490 and BAL1778
and nearly neutral over the large-BAL cohort.

C1 and C5 compose strongly on 1DSfM.  The combination improves every scene,
reduces geometric mean SSE by $21.3\%$, and repairs both C1 losses.  For
example, the Trafalgar ratio changes from $1.0451$ with C1 to $0.7131$ with
C1+C5; NYC Library changes from $1.0299$ to $0.6344$.  This quality costs
$2.54\times$ worker CPU.  On large BAL, C5 does not provide a useful cumulative
gain and slightly worsens BAL3068.  We consequently retain C1+C5 as a 1DSfM
quality configuration, not as the universal practical preset.

These breadth rows are retained pre-recovery artifacts.  We additionally ran a
fresh executable confirmation of the current source over all 15 1DSfM and all
29 BAL scenes.  Table~\ref{tab:stage-c-fresh-cross-family} compares the fresh
global $2\times2$ factorial with a fresh plain DRS control.

\begin{table}[t]
  \centering
  \caption{Fresh current-source K24/I30 cross-family factorial.  SSE and
  optimization ratios use the matched plain DRS row in each family.}
  \label{tab:stage-c-fresh-cross-family}
  \small
  \begin{tabular}{lrrrrrr}
    \hline
    Method & 1DSfM SSE & W/T/L & Opt. & BAL SSE & W/T/L & Opt. \\
    \hline
    C1 & 0.865147 & 13/0/2 & 1.557830 & 0.983780 & 28/0/1 & 1.820914 \\
    C5 & 1.015807 & 5/0/10 & 1.207142 & 1.000057 & 4/15/10 & 1.158716 \\
    C1+C5 & 0.870061 & 13/0/2 & 1.801873 & 0.983638 & 28/0/1 & 2.030313 \\
    \hline
  \end{tabular}
\end{table}

All 176 DRS cases complete.  C1 is the accepted first cumulative rung.  Current
C5 alone is not an aggregate improvement, and adding it to C1 gives
$1.005680\times$ C1 SSE on 1DSfM and $0.999855\times$ on BAL.  We retain both
innovations and the C1+C5 architecture.  This factorial motivated the frozen
global C5 tuning reported below; it was not used to select a mode by scene.  At
this I30 budget every DRS row remains behind Ceres on 1DSfM and approximately
tied on BAL; cross-solver timing is contextual because setup boundaries differ.

A matched $2\times2$ diagnosis shows that the losses are not failed fallback
events.  On Roman, C1 and C5 individually reach $1.0982\times$ and
$1.0556\times$ plain SSE, but their beneficial interaction limits C1+C5 to
$1.0473\times$.  On Yorkminster, the same main effects are $1.0310\times$ and
$1.0637\times$ and interact adversely, producing $1.1843\times$.  BAL135 is a
pure C1 loss of only $0.27\%$ because C5 never activates.  A prefix-identical
Yorkminster continuation with the I30 safeguard schedule held fixed reaches
$0.9080\times$ for C1 and $1.0292\times$ for C1+C5 at I60.  Thus C1's loss is a
finite-horizon lag, while current C5 parameters and their Yorkminster
interaction remain a global tuning diagnostic rather than a scene veto.
The safeguard ensures admissible progress on the current trajectory; it does
not maintain and dominate a shadow plain-DRS trajectory.

We therefore tuned C5 globally and one factor at a time.  Scalar threshold,
window, dwell, and maximum-depth changes did not improve the C1 rung on both
development families.  Delaying C5 until outer iteration five did.  The frozen
common policy uses high/low thresholds $0.35/0.20$, window and dwell three, and
maximum depth two.  It was selected on six development 1DSfM scenes plus five
BAL sentinels, then evaluated unchanged on nine held-out 1DSfM scenes and all
29 BAL scenes.  Tuned C1+C5/C1 geometric SSE is $0.980053$ on all-15 1DSfM and
$0.999486$ on all-29 BAL; tuned C1+C5/plain is $0.847890$ and $0.983275$, with
W/T/L $14/0/1$ and $28/0/1$.  We promote this single global C1+C5 policy while
retaining C1 and C5 as independent ablation switches.

After one excluded warm-up, three measured repeats cover Roman, Trafalgar, and
BAL scenes 52 and 3068.  Endpoint SSE, safeguard counts, and oracle counts are
identical.  Optimization-time coefficients of variation range from
$1.1\%$ to $2.4\%$, and overall-time coefficients range from $0.9\%$ to
$2.5\%$.

We also implement an identical observation-level Huber objective with pixel
scale $0.5$ in the worker, coordinator safeguards, state evaluator, and
left-$\mathrm{SE}(3)$ Ceres reference.  Explicit IRLS weights provide a
positive-semidefinite custom Gauss--Newton model, while exact robust costs drive
trust acceptance.  On Roman/Trafalgar and BAL52/BAL3068, raw-start Huber I30 is
$7.8002\times$ and $1.2674\times$ Ceres geometrically.  Adding 30 Huber
iterations after the tuned L2 state improves that state by only $0.21\%$ and
$0.99\%$, respectively, and remains $2.5698\times$/$1.1679\times$ Ceres.  We
therefore retain robust-objective support but do not promote either current
Huber schedule.

Finally, a frozen cluster-count sweep at $K\in\{2,4,8,16,24\}$ selects two
nondominated global operating points, which we confirm unchanged on all 15
1DSfM and all 29 BAL scenes.  Relative to $K=4$, $K=16$ reaches geometric SSE
$0.953599$ on 1DSfM and $1.002768$ on BAL while reducing optimization time to
$0.584$ and $0.510$.  Its worker CPU rises to $1.456$/$1.568$ and transport to
$1.927$/$1.985$.  We therefore retain $K=4$ as the resource endpoint and
$K=16$ as the latency endpoint, without scene-specific routing.

After one excluded warm-up, three measured repeats of both endpoints on four
sentinels are SSE-bitwise identical in seven of eight scene--$K$ cases; the
maximum relative SSE spread is $2.956\times10^{-9}$.  Rejection and oracle
counts are identical, and the maximum optimization-time coefficient of
variation is $1.75\%$.  Mean $K=16$/$K=4$ optimization-time ratios are $0.5158$
on 1DSfM and $0.4190$ on BAL, with ratio CV below $0.5\%$.

Table~\ref{tab:stage-c-final-reference} separates the complete cohorts and
restores the best base DRS comparison.  The best K1 row is the BAE-style local
diagnostic at I90/T1; K4 and K16 are the frozen C1+C5 I30 endpoints with one
thread per cluster; Ceres is left-$\mathrm{SE}(3)$ at I90/T16.  The three
DRS+Schur rows use the base-backbone Themelis trajectory followed by safeguarded
distributed global Schur corrections and are explicitly polishing workflows,
not DRS-only gains.  Timings are reported within their named CPU boundary, not
as matched work comparisons.

\begin{table}[t]
  \centering
  \caption{Final complete-cohort comparison.  SSE values are geometric ratios
  to independently evaluated Ceres states; time is aggregate optimizer time.}
  \label{tab:stage-c-final-reference}
  \scriptsize
  \begin{tabular}{llrrrr}
    \hline
    Cohort & Method & SSE/Ceres & SSE/base & W/T/L & Time (s) \\
    \hline
    1DSfM-15 & Ceres & 1.000000 & 0.686334 & 0/15/0 & 219.521 \\
    1DSfM-15 & K1 BAE-style & 0.993529 & 0.681892 & 5/0/10 & 2060.937 \\
    1DSfM-15 & Base K24/I200 & 1.457017 & 1.000000 & 1/0/14 & 572.399 \\
    1DSfM-15 & DRS+Schur fast & 1.411625 & 0.968846 & 1/0/14 & 191.161 \\
    1DSfM-15 & DRS+Schur balanced & 1.272048 & 0.873049 & 1/0/14 & 289.496 \\
    1DSfM-15 & DRS+Schur quality & 1.215246 & 0.834064 & 2/0/13 & 381.521 \\
    1DSfM-15 & K4 resource & 2.052536 & 1.408725 & 0/0/15 & 160.607 \\
    1DSfM-15 & K16 latency & 1.957296 & 1.343358 & 0/0/15 & 93.793 \\
    1DSfM-15 & K4 + terminal correction & 1.612187 & 1.106498
      & 0/0/15 & 178.001 \\
    1DSfM-15 & K16 + terminal correction & 1.523013 & 1.045295
      & 1/0/14 & 109.890 \\
    BAL-29 & Ceres & 1.000000 & 1.001599 & 0/29/0 & 1803.780 \\
    BAL-29 & Base K24/I90 & 0.998404 & 1.000000 & 12/0/17 & 1271.575 \\
    BAL-29 & K4 resource & 1.007309 & 1.008920 & 8/0/21 & 939.466 \\
    BAL-29 & K16 latency & 1.010098 & 1.011713 & 8/0/21 & 479.087 \\
    BAL-29 & K4 + terminal correction & 1.005041 & 1.006648
      & 9/0/20 & 1547.694 \\
    BAL-29 & K16 + terminal correction & 1.007594 & 1.009206
      & 9/0/20 & 1000.459 \\
    \hline
  \end{tabular}
\end{table}

Verified BAE coverage is limited to six 1DSfM scenes.  On that matched inset,
Ceres, BAE-style K1, base K24, K4, K16, BAE-CG, and BAE-Nesterov are
respectively $1.000000$, $0.953562$, $1.304394$, $1.922802$, $1.745880$,
$0.944668$, and $1.091773$ times Ceres.
BAE-CG uses an RTX~5090 and is basin-sensitive on Trafalgar; its 56.170-second
GPU time is not compared as a speedup against CPU timings.  No authoritative
K1 BAL-29, BAE 1DSfM-15, or BAE BAL-29 result is available.

Consequently, C1+C5 remains a valid matched-I30 component improvement
($0.847890$ times plain on 1DSfM and $0.983275$ on BAL), but the shorter K4/K16
runs do not replace best base DRS in endpoint quality.  Their contribution is
the measured latency/resource tradeoff; the remaining research problem is to
retain that speed while recovering the longer-horizon base trajectory.  One
terminal distributed Schur correction improves the raw K4/K16 endpoints by
$0.785461/0.778120$ on 1DSfM and $0.997869/0.997663$ on BAL.  These variants
remain above preserved base quality and are reported as separately timed
correction workflows, not as DRS-only gains.

Distributed Schur polishing changes that endpoint frontier without changing
the DRS-only attribution.  Fast I30+10, balanced I60+10, and quality I60+16
reach $0.968846$, $0.873049$, and $0.834064$ times base-I200 SSE at $0.333964$,
$0.505759$, and $0.666529$ times its optimization time.  An I60+20 ceiling
reaches $0.813073$ times base-I200 and $1.184661$ times Ceres in $430.320$ s;
it is a higher-budget quality reference rather than a named deployment preset.

A follow-up hybrid restores the base local Nesterov, persistent DRS trust,
curvature-$0.4$ recovery/decay, and metric scale 75 while retaining shared-only
product-space semantics.  With C1 fixed, applying proposal damping only to
duplicated cameras reaches $0.992977$ times base I30 on all 15 1DSfM scenes and
$0.997428$ on all 29 BAL scenes, with complete execution.  C5 is neutral or
adverse on this backbone.  The equal-I30 gain does not survive the tested I90
continuations: permanent damping reaches $1.017187$/$1.001950$ times base and
exhausts Tower; proposal-only cutoff at I30 reaches $1.022809$/$1.001478$; and
cutting both proposal damping and C1 reaches $1.136224$/$1.009261$.  We report
the hybrid as positive component compatibility evidence, not as a promoted
long-horizon replacement.

\subsection{K1-derived product-space and guarded proposal carryovers}

We next isolate the two K1-derived additions on a copied mature DRS baseline.
Direct left-$\mathrm{SE}(3)$ tangent assembly and exact tangent-metric
consistency are fixed in every arm.  The $2\times2$ factorial varies only
shared-only product-space semantics and a guarded I60+I90 two-direction Krylov
camera/landmark proposal.  This is not a C1--C5 factorial.

\begin{table}[t]
  \centering
  \caption{Copied-baseline K24/I120 K1-carryover factorial.  Ratios below one
  favor the numerator.  S, P, C, and D denote shared-only, proposal-only,
  combined, and direct arms.  Interaction below one indicates that the
  combined effect is better than the product of the isolated effects.}
  \label{tab:k1-carryover-k24-factorial}
  \scriptsize
  \begin{tabular}{lrrrrrr}
    \hline
    Family & S/D & P/D & C/D & P/S & Interaction & C/Ceres \\
    \hline
    1DSfM-15 & 0.992877 & 0.787326 & \textbf{0.762902}
      & 0.768375 & 0.975930 & 1.169700 \\
    BAL-29 & 1.000177 & 0.999998 & 1.000131
      & 0.999953 & 0.999955 & 1.001181 \\
    \hline
  \end{tabular}
\end{table}

On 1DSfM the combined arm improves the direct baseline on 14 of 15 scenes and
both conditional effects are favorable.  Shared-only alone exhausts recovery
on Tower of London, while the combined arm completes it.  On BAL the proposal
is nearly neutral and slightly repairs shared-only; the residual $0.013\%$
combined cost is attributable to the mandatory shared-only factor.  We retain
shared-only semantics as a correctness invariant and the guarded Krylov move as
the associated K24 quality mechanism.

The same frozen stack was then tested without retuning at the established K4
resource and K16 latency budgets on the six-1DSfM/five-BAL development cohort.
All 22 candidate and control trajectories complete without recovery
exhaustion.  Candidate/direct geometric SSE is $0.714243/0.788556$ on 1DSfM
and $1.001972/0.997714$ on BAL at K4/K16.  However, Gendarmenmarkt/K16 reaches
$1.022191\times$ direct, exceeding the predeclared $1.02\times$ scene-tail
bound.  We therefore close this combined-stack transfer without held-out or
all15/all29 expansion and without K-specific tuning.  This decision does not
demote the separately established Stage-C K4/K16 operating points or their
separately labeled terminal-correction variants.

\subsection{Coordinate and finite-solver interactions}

The remaining components are reported as factorials because their effects are
not independent.  Table~\ref{tab:stage-c-factorials} summarizes the decisive
sentinel interactions.  C2 denotes points-$p_{95}$ plus Jacobi equilibration,
and C4 replaces Schur--PCG with finite Schur--Nesterov while holding trust,
damping, nonlinear depth, and safeguards fixed.

\begin{table}[t]
  \centering
  \caption{Long-horizon sentinel interactions.  Values are endpoint SSE ratios.
  ``Raw'' disables scene normalization and camera scaling; ``scaled'' uses the
  maintained points-$p_{95}$/Jacobi coordinates.}
  \label{tab:stage-c-factorials}
  \small
  \begin{tabular}{llrrrr}
    \hline
    Scene & Coordinates & C1 only & C2 only & C1+C2 & Nesterov/PCG with C1 \\
    \hline
    Roman Forum & Raw & 0.763587 & 1.062746 & 0.718185 & 0.883997 \\
    Roman Forum & Scaled & -- & -- & -- & 1.069439 \\
    Trafalgar & Raw & 0.614318 & 0.997396 & 1.042354 & 1.025774 \\
    Trafalgar & Scaled & -- & -- & -- & 0.804377 \\
    BAL1490 & Raw & 0.987128 & 0.991674 & 0.986136 & 0.998941 \\
    BAL1490 & Scaled & -- & -- & -- & 1.000139 \\
    BAL3068 & Raw & 0.991296 & 1.001286 & 1.003482 & 1.019623 \\
    BAL3068 & Scaled & -- & -- & -- & 1.015382 \\
    \hline
  \end{tabular}
\end{table}

C2 is not an automatic cumulative improvement after C1.  Roman exhibits
positive C1--C2 synergy, whereas on Trafalgar scaling changes a strong raw C1
gain ($0.6143$) into a loss ($1.0424$).  We therefore report C2 as an isolated
coordinate study.  C4 is similarly context-dependent: Nesterov repairs scaled
C1 Trafalgar but worsens raw C1 Trafalgar and scaled C1 Roman.  Moreover, equal
configured tolerance does not imply equal delivered work.  PCG terminates at
relative residuals near $6\text{--}7\times10^{-3}$, whereas Nesterov reaches
approximately $8\times10^{-4}\text{--}1.1\times10^{-3}$ and uses
$1.11\text{--}1.33\times$ worker CPU in these contexts.  We treat C4 as a
finite-work basin ablation and make no universal solver-efficiency claim.

The recovered core path permits a broader C4 decision without changing the
distributed consensus method.  Table~\ref{tab:stage-c-c4-all15} reports the
all-15 K24/I30 inner-only gate with outer acceleration disabled.  Coordinates,
damping, preconditioning, trust, local nonlinear depth, and the proximal metric
are fixed; only the finite inner solver changes.

\begin{table}[t]
  \centering
  \caption{Post-crash all-15 C4 gate.  Ratios compare Schur--PCG with finite
  Schur--Nesterov in the matched curvature-recovery family.}
  \label{tab:stage-c-c4-all15}
  \small
  \begin{tabular}{lrrrrr}
    \hline
    Candidate/reference & Summed SSE & Geomean SSE & W/L & Geomean opt. & Worst \\
    \hline
    Schur--PCG/Nesterov & 0.985199 & 1.005578 & 8/7 & 0.931664 & 1.195314 \\
    \hline
  \end{tabular}
\end{table}

PCG is faster on 12 of 15 scenes and improves summed SSE, but its geometric
per-scene SSE and worst-scene tail regress; NYC Library reaches $1.1953\times$
and Piccadilly $1.0913\times$ the Nesterov SSE.  We therefore retain Nesterov
as the named unchanged trajectory and expose PCG only as an independently
switchable, speed-oriented C4 factorial level.  We neither tune the tolerance
from this matrix nor choose the solver by scene.

\subsubsection*{Reproducibility}
Per-scene configurations, artifact paths, endpoint metrics, proximal oracle
counts, worker CPU, worker RSS, repository revision, and the source-recovery
boundary are recorded in the machine-readable supplementary manifest.